\documentclass[10pt,a4paper]{amsart}
\usepackage[margin=19mm]{geometry}
\usepackage{amsmath,amssymb,mathtools}
\usepackage[scr=rsfs,cal=euler]{mathalfa}
\usepackage{xfrac}
\usepackage[unicode,hidelinks,bookmarks=false]{hyperref}
\newcommand{\ie}{i.e.\ }
\newcommand{\cf}{cf.\ }
\newcommand{\resp}{respectively\ }
\newcommand{\Subsection}{\subsection}
\providecommand{\nobreakdash}{\nobreak\hbox{-}}
\setlength{\emergencystretch}{2em}
\multlinegap0pt
\usepackage{mathtools}
\newtheorem{theorem}{Theorem}[section]
\newtheorem{lemma}[theorem]{Lemma}
\newtheorem{corollary}[theorem]{Corollary}
\newtheorem{remark}[theorem]{Remark}
\newcommand{\F}{\mathbb F}
\newcommand{\one}{\mathbf 1}
\newcommand{\AG}{\operatorname{AG}}
\hypersetup{
  pdftitle={Four Special Directions in AG(2,13): The 52-Point Obstruction and the Sharp Minimum},
  pdfauthor={Qihang Wang},
  pdfsubject={Finite affine geometry and special directions},
  pdfkeywords={special direction, equidistributed direction, finite affine plane, polynomial method}
}
\begin{document}
\title[Four Special Directions in AG(2,13): The 52-Point Obstructio]{Four Special Directions in AG(2,13): The 52-Point Obstruction and the Sharp Minimum}
\author{Qihang Wang}
\date{}
\maketitle
\noindent\emph{Source and licence.} Four Special Directions in AG(2,13): The 52-Point Obstruction and the Sharp Minimum. Version arXiv:2609.20495v1 (CC BY 4.0). This material is distributed under the Creative Commons Attribution 4.0 International licence, http://creativecommons.org/licenses/by/4.0/, as recorded in the arXiv OAI licence metadata for this exact version and shown on the abstract page https://arxiv.org/abs/2609.20495v1. Full rights notice: licenses/arxiv-2609.20495-CC-BY-4.0.txt.\par\smallskip
\noindent\textit{Editorial scope.} Counted unit: the original \texttt{theorem} environment \texttt{thm:main} at source lines 108--111 together with its complete proof at lines 321--376; the delivered document keeps source lines 64--389 so that every referenced label and every citation resolves inside it. Native editable source is the paper's own concatenated TeX; the AMS wrapper is editorial formatting only. No publisher class, upstream script or unrelated artwork is redistributed.
\section{Introduction}

Let $S$ be a subset of the affine plane $\F_{13}^{2}$.  Every affine
direction partitions the plane into thirteen parallel lines.  Following the
terminology of Ghidelli \cite{GhidelliRichPoor} and Kiss and Somlai
\cite{KissSomlaiSpecialDirections}, call a direction \emph{special} for $S$
when the thirteen line-intersection counts are not all equal.

The subject belongs to the classical direction problem in finite affine
planes, whose polynomial-method lineage runs from R\'edei's lacunary-polynomial
approach \cite{RedeiLacunary} through later finite-geometric developments such
as Sz\H{o}nyi's direction theorem \cite{SzonyiDirections}; see also Somlai's
recent proof of R\'edei's theorem \cite{SomlaiRedei}.  It also has a
coding-theoretic formulation through multisets and small-weight codewords
supported on concurrent lines; see Adriaensen, Sz\H{o}nyi, and Weiner
\cite{AdriaensenSzonyiWeiner}.  The present paper concerns the ordinary,
unweighted subset problem over the prime field $\F_{13}$.  Related recent
work studies highly structured subsets such as the points below a parabola
and their intersection numbers
\cite{AdriaensenWeinerParabola}.

For a prime $p$ and $1\le n\le p$, Ghidelli proved that a set of cardinality
$np$ in $\F_p^2$ is either
contained in the union of $n$ lines---necessarily a union of $n$ parallel
lines, since only then does the union have cardinality $np$---or has at least
\begin{equation}\label{eq:ghidelli-bound}
  \left\lceil\frac{p+n+2}{n+1}\right\rceil
\end{equation}
special directions \cite[Theorem~1.3]{GhidelliRichPoor}.  Kiss and Somlai
subsequently constructed examples with four special directions meeting this
bound for $p=5,7,11$.  For $p=13$ they constructed a $65$-point example and
asked whether a $52$-point example exists
\cite[Section~6]{KissSomlaiSpecialDirections}.  We prove that no such
$52$-point example exists.  Thus, at $p=13$, the lower-bound value $4p=52$ is not
attainable, whereas the $5p=65$ construction is attainable.

The proof uses the degree bound of Kiss and Somlai
\cite[Proposition~3.1]{KissSomlaiSpecialDirections}: with four special
directions, the residue-valued line-count functions have degree at most two.
We include a self-contained proof of this bound in
Lemma~\ref{lem:degree-two}, classify the resulting profiles in
Section~\ref{sec:profiles}, and derive the contradiction from a
quadratic-character congruence.


\begin{theorem}\label{thm:main}
There is no subset $S\subseteq\F_{13}^{2}$ with $|S|=52$ for which exactly
four affine directions are special.
\end{theorem}


Combining this obstruction with the two cited results gives a sharp numerical
answer, not merely a nonexistence statement at one cardinality.

\begin{corollary}\label{cor:minimum}
The minimum cardinality of a subset of $\F_{13}^{2}$ having exactly four
special directions is $65$.
\end{corollary}

The proof has two logically distinct parts.  The incidence identity in
Section~\ref{sec:incidence} and the polynomial reduction in
Section~\ref{sec:degree} apply to every candidate set.  After that reduction
the line-count functions have degree at most two, and
Section~\ref{sec:profiles} classifies the admissible value tables.  The final
obstruction is a quadratic-character congruence with no solution in
$\F_{13}$.  The argument is therefore not an enumeration of a prefix of the
possible point sets.  Corollary~\ref{cor:minimum} is derived in
Section~\ref{sec:consequences}, where we also state the precise boundary of
what is settled.

\section{The incidence identity}\label{sec:incidence}

For a point $P\in\F_{13}^{2}$ and an affine direction $d$, let $c_d(P)$
denote the number of points of $S$ on the line of direction $d$ through $P$.
Write $\one_S(P)=1$ if $P\in S$, and $\one_S(P)=0$ otherwise.

\begin{lemma}\label{lem:incidence}
Suppose that $|S|=52$ and that the set $D$ of special directions has size
four.  Then
\begin{equation}\label{eq:four-counts}
  \sum_{d\in D}c_d(P)=12+13\one_S(P)
  \qquad(P\in\F_{13}^{2}).
\end{equation}
\end{lemma}

\begin{proof}
In a nonspecial direction the thirteen line counts are equal and sum to $52$,
so each count is $4$.  There are fourteen directions in $\AG(2,13)$.  When
we sum $c_d(P)$ over all of them, every selected point $Q\ne P$ is counted
once, while $P$ itself is counted fourteen times if $P\in S$.  Hence
\[
  \sum_d c_d(P)=52+13\one_S(P).
\]
The ten nonspecial directions contribute $40$, which gives
\eqref{eq:four-counts}.
\end{proof}

An invertible affine linear change of coordinates sends two distinct
directions in $D$ to the vertical and horizontal directions.  A residual
scaling of the two axes then sends a third to slope $1$.  The fourth has
slope $r\in\F_{13}\setminus\{0,1\}$.  Thus every four-direction set is
represented as
\[
  D=\{\infty,0,1,r\},
  \qquad r\in\{2,3,\ldots,12\}.
\]
Since the argument applies to every such $r$, no choice of projective-orbit
representatives is needed.

Write $V,H,A,B\colon\F_{13}\to\{0,1,\ldots,13\}$ for the integer line-count
functions in these four directions.  With line labels
\[
  \ell_\infty(x,y)=x,
  \qquad \ell_m(x,y)=y-mx,
\]
reduction of \eqref{eq:four-counts} modulo $13$ gives
\begin{equation}\label{eq:polynomial-identity}
  V(x)+H(y)+A(y-x)+B(y-rx)=12
  \qquad((x,y)\in\F_{13}^{2}).
\end{equation}
Here and below the same letters denote the residue-valued functions and their
unique polynomial representatives of degree at most $12$.

\section{Reduction to quadratic line-count functions}\label{sec:degree}

\begin{lemma}\label{lem:rank}
For every $r\in\F_{13}\setminus\{0,1\}$ and every
$k\in\{3,4,\ldots,12\}$, the four homogeneous forms
\[
  x^k,\qquad y^k,\qquad (y-x)^k,\qquad (y-rx)^k
\]
are linearly independent over $\F_{13}$.
\end{lemma}

\begin{proof}
A homogeneous polynomial of degree $k$ is determined by its restriction to
the line $y=1$, so it is equivalent to show that the dehomogenizations
\[
  x^k,\qquad 1,\qquad (1-x)^k,\qquad (1-rx)^k
\]
are linearly independent in $\F_{13}[x]$.  Extracting the coefficients of
$1$, $x$, $x^{2}$, and $x^{k}$ produces the $4\times4$ matrix
\[
\begin{pmatrix}
0 & 1 & 1 & 1 \\
0 & 0 & -k & -kr \\
0 & 0 & \binom{k}{2} & \binom{k}{2}r^{2} \\
1 & 0 & (-1)^{k} & (-r)^{k}
\end{pmatrix}.
\]
Its determinant is $-k\binom{k}{2}r(1-r)$.  For
$k\in\{3,\ldots,12\}$ and $r\in\F_{13}\setminus\{0,1\}$ each factor is
nonzero in $\F_{13}$, so the four forms are linearly independent.
\end{proof}

\begin{lemma}\label{lem:degree-two}
Each of $V,H,A,B$ has degree at most two.
\end{lemma}

\begin{proof}
Interpret $V,H,A,B$ as their unique polynomial representatives of degree at
most $12$.  The left-hand side of \eqref{eq:polynomial-identity}, minus
$12$, is then a bivariate polynomial of degree at most $12$ in each
variable, agreeing as a function with a quantity that vanishes on the full
$13\times13$ grid.  For each fixed $y_{0}\in\F_{13}$ the resulting univariate
polynomial in $x$ has $13$ roots and degree less than $13$, hence is the
zero polynomial; the coefficient polynomials in $y$ likewise vanish, so the
bivariate polynomial is identically zero.

If one of the four univariate polynomials had degree at least three, let $k$
be the highest such degree.  The homogeneous degree-$k$ part of the identity
would be a nontrivial linear relation among the four forms in
Lemma~\ref{lem:rank}.  This contradiction proves the assertion.
\end{proof}

\begin{remark}[Attribution]\label{rem:attribution}
Kiss and Somlai proved that if a
subset of $\F_p^{2}$ of cardinality divisible by $p$ has $k\ge2$ special
directions, then the mod-$p$ line-count function of every direction is a
polynomial of degree at most $k-2$
\cite[Proposition~3.1]{KissSomlaiSpecialDirections}; the statement above is
the case $k=4$, $p=13$, restricted to the four special directions.  The
short proof via Lemma~\ref{lem:rank} is included to keep the paper
self-contained.
\end{remark}

\section{Exact classification of the reduced profiles}\label{sec:profiles}

Let $\chi$ denote the quadratic character of $\F_{13}^{\times}$.

\begin{lemma}\label{lem:profiles}
Let $q\in\F_{13}[t]$ have degree at most two.  Suppose that its thirteen
values lift to integers in $\{0,1,\ldots,13\}$ whose sum is $52$.
Then exactly one of the following occurs.
\begin{enumerate}
\item The residue polynomial is the constant $4$.
\item The residue polynomial is the constant $0$, with exactly four of the
      thirteen values lifted to $13$.
\item The polynomial is genuinely quadratic,
      \[
        q(t)=at^2+bt+c,
        \qquad a\ne0,
      \]
      no zero value is lifted to $13$, and
      \begin{equation}\label{eq:vertex-condition}
        c-\frac{b^2}{4a}=6-2\chi(a).
      \end{equation}
\end{enumerate}
No nonconstant linear polynomial is admissible.
\end{lemma}

\begin{proof}
An integer line count is between $0$ and $13$.  Once its residue modulo $13$
is nonzero, its lift is unique; a zero residue may lift to either $0$ or $13$.
For a constant residue $c\neq0$ the sum of lifts is $13c$, so $c=4$.  For the
zero polynomial the sum is $13$ times the number of $13$-lifts, so exactly
four lines are full.

A nonconstant linear polynomial permutes $\F_{13}$.  The sum of its least
nonnegative residues is $0+1+\cdots+12=78$; lifting its unique zero changes
this only to $91$.  Neither value is $52$.

Now suppose $q(t)=at^{2}+bt+c$ with $a\neq0$.  Completing the square gives
\[
  q(t)=a\Bigl(t+\frac{b}{2a}\Bigr)^{2}+\delta,
  \qquad
  \delta=c-\frac{b^{2}}{4a}.
\]
As $t$ runs through $\F_{13}$, so does the translated argument, and the value
multiset is $\{\delta+au^{2}:u\in\F_{13}\}$.  Put
\[
 Q_+=\{1,3,4,9,10,12\},\qquad Q_-=\{2,5,6,7,8,11\}.
\]
If $\varepsilon=\chi(a)$, the vertex value $\delta$ occurs once, and
each value $\delta+v$ with $v\in Q_\varepsilon$ occurs twice.
For $z\in\F_{13}$ let $[z]$ denote its representative in
$\{0,\ldots,12\}$.  The sum of the least nonnegative residues is therefore
\[
 T_\varepsilon(\delta)=[\delta]+2\sum_{v\in Q_\varepsilon}[\delta+v].
\]
The complete table is
\[
\setlength{\arraycolsep}{3.5pt}
\begin{array}{c|rrrrrrrrrrrrr}
\delta&0&1&2&3&4&5&6&7&8&9&10&11&12\\ \hline
T_+(\delta)&78&65&78&65&52&65&78&91&104&91&78&91&78\\
T_-(\delta)&78&91&78&91&104&91&78&65&52&65&78&65&78
\end{array}
\]
Lifting any zero to $13$ only increases the sum.  Every entry is at least
$52$, with equality precisely at $(\varepsilon,\delta)=(1,4)$ or
$(-1,8)$.  Thus these two cases, without any zero lift, are exactly the
possibilities in \eqref{eq:vertex-condition}.  For each of the $12$ choices
of $a\ne0$ and $13$ choices of $b$, that condition determines $c$ uniquely;
hence exactly $156$ of the $2028$ quadratic coefficient triples are admissible.
\end{proof}

\section{The character obstruction}


\begin{proof}[Proof of Theorem~\ref{thm:main}]
Assume that such a set $S$ exists, and use the normalization above.  Compare
the homogeneous quadratic terms of \eqref{eq:polynomial-identity}.  If
$\lambda$ is the coefficient of $t^2$ in $B$, then the quadratic coefficients
of $(V,H,A,B)$ are
\begin{equation}\label{eq:quadratic-coefficients}
  \lambda\bigl(-r(r-1),\ r-1,\ -r,\ 1\bigr).
\end{equation}

If $\lambda=0$, all four count polynomials have degree at most one.  Each of
the four directions is special, so Lemma~\ref{lem:profiles} rules out the
constant-$4$ and nonconstant-linear cases.  Every polynomial must therefore
be the constant residue zero, making the left-hand side of
\eqref{eq:polynomial-identity} congruent to zero rather than $12$.  This is a
contradiction.

Suppose now that $\lambda\ne0$.  If $u$ and $w$ are the linear coefficients
of $A$ and $B$, comparison of the degree-one terms gives the linear
coefficients of $(V,H,A,B)$ as
\begin{equation}\label{eq:linear-coefficients}
  (u+rw,\ -u-w,\ u,\ w).
\end{equation}
All four leading coefficients in \eqref{eq:quadratic-coefficients} are
nonzero, so Lemma~\ref{lem:profiles} applies in its quadratic case.  For a
quadratic $at^2+bt+c$, solve \eqref{eq:vertex-condition} for the constant
term:
\[
  c=6-2\chi(a)+\frac{b^2}{4a}.
\]

The sum of the four square-completion corrections vanishes.  Indeed, since
$4\lambda\neq0$ in $\F_{13}$, it is equivalent to multiply through by
$4\lambda$; the coefficients of $u^{2}$, $uw$, and $w^{2}$ are then
respectively
\begin{align*}
 &\frac1{-r(r-1)}+\frac1{r-1}+\frac1{-r}=0,\\
 &\frac{2r}{-r(r-1)}+\frac2{r-1}=0,\\
 &\frac{r^2}{-r(r-1)}+\frac1{r-1}+1=0.
\end{align*}
Consequently the constant term in \eqref{eq:polynomial-identity} would require
\begin{equation}\label{eq:character-equation}
  24-2\chi(\lambda)\Sigma\equiv12\pmod{13},
\end{equation}
where
\[
  \Sigma=
  \chi\bigl(-r(r-1)\bigr)+\chi(r-1)+\chi(-r)+1.
\]
The integer $\Sigma$ is a sum of four values $\pm1$, hence lies in
$\{-4,-2,0,2,4\}$.  Reducing \eqref{eq:character-equation} gives
$\chi(\lambda)\Sigma\equiv6\pmod{13}$, so the congruence asks for
$\Sigma\equiv6\pmod{13}$ when $\chi(\lambda)=1$, or
$\Sigma\equiv-6\equiv7\pmod{13}$ when $\chi(\lambda)=-1$.  Neither residue
occurs among the five displayed integers.  This final contradiction proves
the theorem.
\end{proof}


\begin{remark}[Scope of the finite classification]
The only case-by-case arithmetic is the classification of quadratic value
tables in Lemma~\ref{lem:profiles}, equivalently the $2028$ coefficient
triples $(a,b,c)$.  The rank identity of Lemma~\ref{lem:rank} and the
character congruence \eqref{eq:character-equation} are closed-form and
independent of a choice of $r$.  Neither step enumerates $52$-point subsets:
the incidence identity and polynomial reduction are what make the quadratic
classification exhaustive for all such subsets and all four-direction sets.
\end{remark}

\section{The sharp minimum and the remaining scope}\label{sec:consequences}


\begin{thebibliography}{99}
\bibitem[AdriaensenSzonyiWeiner]{AdriaensenSzonyiWeiner} S.~Adriaensen, T.~Sz\H{o}nyi, and Z.~Weiner, \emph{Multisets with few special directions and small weight codewords in Desarguesian planes}, Des.\ Codes Cryptogr. \textbf{94} (2026), no.~2, article 35, \href{https://doi.org/10.1007/s10623-025-01777-8} {doi:10.1007/s10623-025-01777-8}.
\bibitem[AdriaensenWeinerParabola]{AdriaensenWeinerParabola} S.~Adriaensen and Z.~Weiner, \emph{Points below a parabola in affine planes of prime order}, Bull. Belg. Math. Soc. Simon Stevin \textbf{32} (2025), no.~3, 332--342, \href{https://doi.org/10.36045/j.bbms.241203} {doi:10.36045/j.bbms.241203}.
\bibitem[GhidelliRichPoor]{GhidelliRichPoor} L.~Ghidelli, \emph{On rich and poor directions determined by a subset of a finite plane}, Discrete Math. \textbf{343} (2020), no.~5, article 111811, \href{https://doi.org/10.1016/j.disc.2020.111811} {doi:10.1016/j.disc.2020.111811}.
\bibitem[KissSomlaiSpecialDirections]{KissSomlaiSpecialDirections} G.~Kiss and G.~Somlai, \emph{Special directions on the finite affine plane}, Des.\ Codes Cryptogr. \textbf{92} (2024), no.~9, 2587--2597, \href{https://doi.org/10.1007/s10623-024-01404-y} {doi:10.1007/s10623-024-01404-y}.
\bibitem[RedeiLacunary]{RedeiLacunary} L.~R\'edei, \emph{Lacunary polynomials over finite fields}, North-Holland Publishing Co., Amsterdam, 1973. Translated from the German by I.~F\"oldes.
\bibitem[SomlaiRedei]{SomlaiRedei} G.~Somlai, \emph{A new proof of R\'edei's theorem on the number of directions}, Arch. Math. (Basel) \textbf{122} (2024), no.~6, 575--580, \href{https://doi.org/10.1007/s00013-024-01979-x} {doi:10.1007/s00013-024-01979-x}.
\bibitem[SzonyiDirections]{SzonyiDirections} T.~Sz\H{o}nyi, \emph{On the number of directions determined by a set of points in an affine Galois plane}, J. Combin. Theory Ser. A \textbf{74} (1996), no.~1, 141--146, \href{https://doi.org/10.1006/jcta.1996.0042} {doi:10.1006/jcta.1996.0042}.
\end{thebibliography}
\end{document}
